% =============================================================================
%  s03_conclusiones.tex — Sección 3: Conclusiones y recomendaciones
% =============================================================================

\section{Conclusiones}

% ── Frame: Síntesis ───────────────────────────────────────────────────────────
\begin{frame}{Síntesis de Hallazgos: Diagnóstico y Estructura}
  \begin{columns}[T]
    \begin{column}{0.48\textwidth}
      \begin{block}{Lo que funciona (API)}
        \begin{itemize}
          \item El \termino{Anillo Periférico Interior (API)} actúa como una frontera técnica efectiva 
                entre zonas consolidadas y periferias.
          \item Los corredores circulares funcionan como motores de recalibración para \termino{coser} 
                retículas desfasadas.
          \item La redundancia anillar reduce la vulnerabilidad sistémica ante fallos en un 53\%.
        \end{itemize}
      \end{block}
    \end{column}
    \begin{column}{0.48\textwidth}
      \begin{alertblock}{Los retos pendientes (APE)}
        \begin{itemize}
          \item Vulnerabilidad crítica en \textit{Macro-Hubs} radiales por falta de caminos alternos transversales.
          \item El término \termino{Periférico} es insuficiente si no engloba el crecimiento hacia el cuarto contorno.
          \item Inexistencia de un esquema funcional para el \termino{Anillo Periférico Exterior (APE)}.
        \end{itemize}
      \end{alertblock}
    \end{column}
  \end{columns}
\end{frame}

% ── Frame: Recomendaciones ───────────────────────────────────────────────────
\begin{frame}{Lineamientos de Política Pública y Gobernanza}
  \begin{enumerate}
    \item \textbf{Redefinición Regional (API vs. APE):} \\
          Acotar el API a la escala urbana central y formalizar el \termino{Anillo Periférico Exterior (APE)} como 
          eje de integración metropolitana regional.
    % \pause
    \item \textbf{Gobernanza Metropolitana:} \\
          Crear un organismo con atribuciones vinculantes sobre CDMX, EDOMEX, Morelos e Hidalgo para la gestión del 
          APE.
    % \pause
    \item \textbf{Reducción de Vulnerabilidad Operativa:} \\
          Priorizar inversiones en infraestructura transversal sobre el trazo de fuerzas capilares identificadas 
          para descentralizar los \textit{hubs} históricos.
    % \pause
    \item \textbf{Recalibración mediante VFTModel:} \\
          Vincular los planes de desarrollo a indicadores de \textit{Robustez Geométrica ($\Delta E$)} y 
          conectividad lateral en el tercer y cuarto contorno.
  \end{enumerate}
\end{frame}

% ── Frame: Cierre ────────────────────────────────────────────────────────────
\begin{frame}[plain]
  \begin{center}
    \vfill
    {\color{ColorAcento}\rule{0.6\textwidth}{1.5pt}}\\[0.5cm]
    {\Large\bfseries\color{ColorPrincipal} Gracias}\\[0.3cm]
    {\normalsize \Autor}\\[0.1cm]
    {\small \MateriaCorta}\\[0.1cm]
    {\small \Semestre}\\[0.5cm]
    {\color{ColorAcento}\rule{0.6\textwidth}{1.5pt}}
    \vfill
  \end{center}
\end{frame}